% JETSPIN running documentation

\section{Running JETSPIN\label{sec:Running-JETSPIN}}

In order to run the JETSPIN executable (main.x) the user should prepare
the input file following the hints provided in Section \ref{sec:Input-file}.
The test cases located in \textit{examples} sub-directory can be used
as starting point for editing the input file. As JETSPIN is executed,
the program will look for the input file, named \texttt{input.dat},
in the same directory, where the executable was called. For instance,
if JETSPIN is called in the \textit{execute} sub-directory, the code
will look for the file \texttt{input.dat}, which must have been put
in the same sub-directory.

The JETSPIN code can be executed in serial mode by the command\texttt{ }

\bigskip{}

\texttt{./main.x}

\bigskip{}

or in parallel mode by the command

\bigskip{}

\texttt{mpirun -np \$ncpu ./main.x}

\bigskip{}
 where \texttt{\$ncpu} is the number of CPUs used in parallel (e.g.
\texttt{mpirun -np 4 ./main.x} in order to run JETSPIN in parallel
mode on 4 CPUs).

A successful run of JETSPIN will generate several output files, which
appear in the directory, where the executable was called. The output
files will be produced only if their writing was explicitly requested
in input file by using proper directives (see Section \ref{sec:Input-file}).
A detailed description of the generated output files is provided in
Section \ref{sec:Output-files}. Briefly, the main output files are
the files \texttt{statout.dat},\texttt{ traj.xyz} and \texttt{frame\%06d.xyz},
which are human readable. The \texttt{statout.dat} file contains a catalogue of instantaneous
or blocked averaged values of several variables, which can be used
for temporal or statistical analysis in post-processing. The \texttt{traj.xyz}
file provides a time ordered sequence of configurations to facilitate
further analysis of the jet dynamics, while the \texttt{frame\%06d.xyz}
file, reports the same information in splitted files, one for each
registered time frame. Further, as JETSPIN is running, a series of
lines will be printed in the terminal window, providing several information
as requested by suitable directives in input file (see Section \ref{sec:Input-file}).
Also present will be the binary \texttt{save.dat} file, which is
\textit{not} human readable. It is written every \textit{restart dump}
steps (100,000 by default) and at the end of the run, and holds the bead
state in double precision followed by the state of the run: the position in
the Gaussian pool, the state of the random-number generator, the counters of
the dynamic refinement, and whether the device step of the OpenACC build has
engaged. Together with \texttt{input.dat} it restarts the job (see Section
\ref{sec:Restarting-JETSPIN}).

\subsection{Running on a GPU\label{sec:Running-GPU}}

The executable of the \texttt{nvfortran-openacc} target
(Sec.~\ref{sec:Compiling-the-Source}) is a serial program that uses one
GPU; it is started as \texttt{./main.x}. It moves the whole time step to
the GPU (the device step) for the Euler, RK2, and RK4 integrators on
\texttt{system 3}, Maxwell or Kelvin--Voigt, with or without evaporation and
air drag, without dynamic refinement (or the other options that tag beads),
and for the Platen integrator on \texttt{system 4}, Maxwell, with or
without evaporation, insertion, and dynamic refinement. Multiple-step
Coulomb sums, time-dependent external fields, the drag velocity, bead
tracking, fixed bead sets with \textit{removing yes}, one-dimensional
systems, MPI runs, and the developer options Lorentz force, upper potential
(\textit{wall}), variable mass and breakup are not supported yet.

The device step engages when \texttt{npjet}, the index of the nozzle bead,
reaches 100 (it counts the beads in the arrays, collected ones included,
not only the active ones), and stays engaged when a compaction lowers it
and across a restart, since \texttt{save.dat} records it. The terminal then
shows \texttt{OpenACC device step engaged (scheme, model) at step N with M
active beads}, or \texttt{OpenACC device step resumed (...)} followed by
\texttt{engaged before the restart} in a restarted run. Below that point
the executable runs the code of the CPU build and reproduces, byte for byte,
the results of a CPU build of the same compiler version. Above it the
results agree with the CPU build within roundoff: the device sums
accumulate in another order, so trajectories that amplify roundoff
separate. The runs that the device step does not support execute the CPU
code too, with only the direct Coulomb sums offloaded once the jet has at
least 128 active beads. A device stage that cannot evaluate the equations
of motion of a run stops it with error 22, a guard that a supported run
does not reach.

The following environment variables, set to \texttt{1} unless stated
otherwise, change this behaviour:
\begin{itemize}
\item \texttt{JETSPIN\_OPENACC\_COULOMB\_MIN\_BEADS=}$n$ sets the 128-bead
threshold of the offloaded Coulomb sums ($0$ offloads every call);
\item \texttt{JETSPIN\_OPENACC\_DISABLE\_PERSISTENT} or
\texttt{JETSPIN\_OPENACC\_DISABLE\_EOM} keeps the device step closed, so
that the run executes the CPU code with the offloaded Coulomb sums above;
\item \texttt{JETSPIN\_OPENACC\_DISABLE\_COULOMB} keeps the one-dimensional
Coulomb sum on the host;
\item \texttt{JETSPIN\_OPENACC\_SYNC} makes the Platen device step of a run
with insertion synchronous (the other device steps always are);
\item \texttt{JETSPIN\_PROFILE} (also \texttt{yes}, \texttt{true},
\texttt{on}) prints the wall time of the main parts of the time loop at the
end of the run;
\item \texttt{JETSPIN\_TOPOLOGY\_SNAPSHOT} writes the active state at
every insertion and removal to \texttt{topology-state.dat}.
\end{itemize}
The last two act in the CPU builds too. Two development variables, which
set the spare capacity of refining runs, are described in
\texttt{docs/introduction/compiling.md}.
For timing, the process should be bound to the NUMA node of its GPU
(\texttt{nvidia-smi topo -m} lists it), for instance with
\texttt{numactl --cpunodebind=N --membind=N ./main.x}: the placement of
the host process alone changes the time of latency-bound runs by up to
8\%.
